\documentclass[11pt,a4paper]{article}
\usepackage{fontspec}
\usepackage{xeCJK}
\setCJKmainfont{AppleGothic}
\setmainfont{Menlo}
\usepackage[margin=2cm]{geometry}
\usepackage{listings}
\usepackage{xcolor}
\usepackage{hyperref}
\usepackage{fancyhdr}
\usepackage{titlesec}
\usepackage{enumitem}
\usepackage{booktabs}
\usepackage{longtable}
\usepackage{tcolorbox}

\definecolor{codebg}{HTML}{F5F5F5}
\definecolor{codeframe}{HTML}{CCCCCC}
\definecolor{keyword}{HTML}{0000FF}
\definecolor{comment}{HTML}{008000}
\definecolor{string}{HTML}{A31515}

\lstset{
  basicstyle=\small\ttfamily,
  backgroundcolor=\color{codebg},
  frame=single,
  rulecolor=\color{codeframe},
  breaklines=true,
  breakatwhitespace=false,
  tabsize=2,
  showstringspaces=false,
  keywordstyle=\color{keyword},
  commentstyle=\color{comment},
  stringstyle=\color{string},
  xleftmargin=0.5em,
  xrightmargin=0.5em,
  aboveskip=0.5em,
  belowskip=0.5em,
}

\lstdefinelanguage{cmake}{
  morekeywords={cmake_minimum_required,project,set,add_library,add_executable,target_link_libraries,find_package,include,list,include_directories,target_include_directories,add_compile_options,APPEND,PUBLIC,PRIVATE,REQUIRED,CONFIG,VERSION,LANGUAGES},
  morecomment=[l]{\#},
}

\pagestyle{fancy}
\fancyhf{}
\fancyhead[L]{\leftmark}
\fancyhead[R]{\thepage}

\titleformat{\section}{\Large\bfseries}{}{0em}{}[\titlerule]
\titleformat{\subsection}{\large\bfseries}{}{0em}{}

\title{\Huge\bfseries C++ / MLIR 학습 자료 통합본\\[0.5em]\large Mini KV Store + MLIR Compiler Study}
\author{kv-cli project}
\date{\today}

\begin{document}
\maketitle
\tableofcontents
\newpage

% ============================================================
\section{프로젝트 개요 (README)}
% ============================================================

터미널에서 동작하는 작은 key-value 저장소를 C++17로 만든다. 데이터는 파일에 저장되어 프로그램을 다시 실행해도 유지된다.

\subsection{사용법}
\begin{lstlisting}[language=bash]
./kvstore set name alice    # store a key-value pair
./kvstore get name          # retrieve a value
./kvstore delete name       # remove a key
./kvstore list              # show all pairs (sorted)
./kvstore count             # number of stored keys
./kvstore export json       # export as JSON
./kvstore export csv        # export as CSV
\end{lstlisting}

\subsection{프로젝트 구조}
\begin{lstlisting}
practice/mini-kvstore/
|- CMakeLists.txt
|- include/
|   |- store.h        # KVStore class
|   |- command.h      # Command interface + concrete commands
|   +- exporter.h     # Exporter interface + JSON/CSV
|- src/
|   |- main.cpp       # CLI parsing + command dispatch
|   |- store.cpp      # KVStore implementation
|   |- command.cpp    # Command implementations
|   +- exporter.cpp   # Exporter implementations
+- test.sh            # Manual test script
\end{lstlisting}

\subsection{연습하는 C++ 개념}
\begin{itemize}[nosep]
\item \textbf{Command Pattern} --- encapsulate each CLI action as an object
\item \textbf{Factory Pattern} --- \texttt{parseCommand()}, \texttt{createExporter()}
\item \textbf{Strategy Pattern} --- swap export formats via polymorphism
\item \textbf{Virtual Interface} --- \texttt{Command} and \texttt{Exporter} base classes
\item \textbf{RAII} --- file I/O with \texttt{ifstream}/\texttt{ofstream}
\item \textbf{std::optional} --- safe "not found" returns
\item \textbf{std::unique\_ptr} --- ownership without manual \texttt{delete}
\item \textbf{std::map} --- sorted associative container (red-black tree)
\item \textbf{Header/implementation separation} --- \texttt{.h} declarations, \texttt{.cpp} definitions
\end{itemize}

\newpage

% ============================================================
\section{1장. LLVM/MLIR 빌드 시스템}
% ============================================================

\subsection{LLVM 빌드의 전체 흐름}

\begin{lstlisting}
소스코드 다운로드 -> CMake 설정(configure) -> 빌드(compile) -> 설치(install)
\end{lstlisting}

\textbf{Step 1: 소스 받기}
\begin{lstlisting}[language=bash]
git clone https://github.com/llvm/llvm-project.git
cd llvm-project
git checkout llvmorg-18.1.0   # 특정 버전 고정 (중요!)
\end{lstlisting}

\textbf{Step 2: CMake 설정}
\begin{lstlisting}[language=bash]
cmake -S llvm -B build \
  -DCMAKE_BUILD_TYPE=Release \
  -DCMAKE_INSTALL_PREFIX=$HOME/llvm-install \
  -DLLVM_ENABLE_PROJECTS=mlir \
  -DLLVM_TARGETS_TO_BUILD=host \
  -DLLVM_ENABLE_ASSERTIONS=ON
\end{lstlisting}

\begin{longtable}{p{5cm}p{9cm}}
\toprule
옵션 & 의미 \\
\midrule
\texttt{-S llvm} & 소스 디렉토리 \\
\texttt{-B build} & 빌드 디렉토리 (out-of-source) \\
\texttt{-DCMAKE\_BUILD\_TYPE=Release} & 최적화 빌드 (Debug는 10GB+) \\
\texttt{-DCMAKE\_INSTALL\_PREFIX} & 설치 경로 \\
\texttt{-DLLVM\_ENABLE\_PROJECTS=mlir} & MLIR도 같이 빌드 \\
\texttt{-DLLVM\_TARGETS\_TO\_BUILD=host} & 현재 CPU만 빌드 \\
\texttt{-DLLVM\_ENABLE\_ASSERTIONS=ON} & assert 활성화 \\
\bottomrule
\end{longtable}

\textbf{Step 3\&4: 빌드 + 설치}
\begin{lstlisting}[language=bash]
cmake --build build -j$(nproc)
cmake --install build
\end{lstlisting}

\subsection{CMake 핵심 명령어 5개}
\begin{lstlisting}[language=cmake]
project(GaweeMLIR LANGUAGES CXX)
add_library(GaweeDialect lib/Gawee/GaweeDialect.cpp)
add_executable(gawee-opt tools/gawee-opt.cpp)
target_link_libraries(gawee-opt PRIVATE GaweeDialect MLIRIR)
find_package(MLIR REQUIRED CONFIG)
\end{lstlisting}

\subsection{find\_package 동작 원리}
\texttt{find\_package(MLIR REQUIRED CONFIG)} 한 줄이 하는 일:
\begin{enumerate}[nosep]
\item \texttt{-DMLIR\_DIR} 경로에서 \texttt{MLIRConfig.cmake}를 찾는다
\item 그 안의 변수들 (\texttt{MLIR\_INCLUDE\_DIRS} 등)을 로드한다
\end{enumerate}

\subsection{PUBLIC vs PRIVATE}
\begin{longtable}{p{3cm}p{5cm}p{5cm}}
\toprule
 & PUBLIC & PRIVATE \\
\midrule
의미 & 나도 쓰고, 나를 쓰는 사람도 씀 & 나만 씀 \\
전파 & 자동 전파됨 & 전파 안 됨 \\
언제 & 라이브러리에 사용 & 최종 실행파일에 사용 \\
\bottomrule
\end{longtable}

\newpage

% ============================================================
\section{2장. C++ 핵심 개념}
% ============================================================

\subsection{Template vs Virtual --- 정적 vs 동적 다형성}

\begin{longtable}{p{3cm}p{5cm}p{5cm}}
\toprule
 & Virtual (동적) & Template (정적) \\
\midrule
결정 시점 & 런타임 (vtable 조회) & 컴파일 타임 (코드 생성) \\
오버헤드 & 함수 포인터 간접 호출 & 없음 (인라인 가능) \\
단점 & 매 호출마다 간접 비용 & 코드 팽창 \\
유연성 & 런타임에 타입 결정 가능 & 컴파일 타임에 확정 \\
\bottomrule
\end{longtable}

\begin{lstlisting}[language=C++]
// Virtual
class Animal {
  virtual void speak() = 0;
};
class Dog : public Animal {
  void speak() override { printf("woof"); }
};

// Template
template <typename T>
void process(T &item) { item.speak(); }
\end{lstlisting}

\subsection{포인터 총정리}

\begin{longtable}{p{3.5cm}p{2.5cm}p{2cm}p{2cm}p{2cm}}
\toprule
종류 & 소유권 & Nullable & 자동 해제 \\
\midrule
\texttt{T*} (raw) & 없음 & O & X \\
\texttt{unique\_ptr<T>} & 단독 소유 & O & O \\
\texttt{shared\_ptr<T>} & 공유 소유 & O & O \\
\texttt{T\&} (reference) & 없음 & X & X \\
\bottomrule
\end{longtable}

\subsection{Move Semantics}

\begin{lstlisting}[language=C++]
// move = 소유권 이전, 원본은 이후 사용 금지
applyPartialConversion(module, target, std::move(patterns));
// patterns는 이 줄 이후 빈 껍데기

// return에서 자동 move
std::unique_ptr<Pass> createMyPass() {
  return std::make_unique<MyPass>();  // 자동 move
}
\end{lstlisting}

\subsection{LLVM 캐스팅 3형제}

\begin{longtable}{p{4cm}p{4cm}p{4cm}}
\toprule
LLVM & Python 대응 & 실패 시 \\
\midrule
\texttt{isa<T>(x)} & \texttt{isinstance(x, T)} & false \\
\texttt{cast<T>(x)} & \texttt{(T)x} & crash \\
\texttt{dyn\_cast<T>(x)} & \texttt{x as T or None} & nullptr \\
\bottomrule
\end{longtable}

\subsection{Value vs Reference Semantics}

\begin{lstlisting}[language=C++]
// C++ 기본 = 값 (복사)
std::vector<int> a = {1, 2, 3};
std::vector<int> b = a;    // 복사! 독립적
b.push_back(4);            // a는 여전히 {1, 2, 3}

// 복사 방지: const T&
void process(const std::vector<int> &data) { ... }
\end{lstlisting}

\newpage

% ============================================================
\section{2b장. C++ 코드 패턴 (Python $\leftrightarrow$ C++)}
% ============================================================

\subsection{패턴 1: CLI (Python argparse $\to$ LLVM cl::opt)}
\begin{lstlisting}[language=C++]
static cl::opt<std::string> inputFilename(
    cl::Positional, cl::desc("<input>"), cl::Required);

int main(int argc, char **argv) {
    InitLLVM y(argc, argv);
    cl::ParseCommandLineOptions(argc, argv, "My tool\n");
    auto fileOrErr = MemoryBuffer::getFileOrSTDIN(inputFilename);
}
\end{lstlisting}

\subsection{패턴 2: 에러 체인 (try/except $\to$ if/return)}
\begin{lstlisting}[language=C++]
auto fileOrErr = MemoryBuffer::getFileOrSTDIN(inputFilename);
if (auto ec = fileOrErr.getError()) {
    errs() << "Error: " << ec.message() << "\n";
    return 1;
}
\end{lstlisting}

\subsection{패턴 3: 처리기 클래스}
\begin{lstlisting}[language=C++]
class MLIREmitter {
public:
    explicit MLIREmitter(MLIRContext *context);
    OwningOpRef<ModuleOp> emit(const json::Object &graph);
private:
    MLIRContext *ctx;                            // 빌려 쓰는 포인터
    std::unique_ptr<OpBuilder> builder;          // 소유
    std::unordered_map<std::string, Value> valueMap;
};
\end{lstlisting}

\subsection{패턴 4: Visitor 디스패치}
\begin{lstlisting}[language=C++]
if (*opType == "Conv")     return emitConv(node, values);
if (*opType == "Relu")     return emitRelu(node, values);
if (*opType == "Add")      return emitAdd(node, values);
\end{lstlisting}

\subsection{패턴 5: Registry (patterns.add)}
\begin{lstlisting}[language=C++]
RewritePatternSet patterns(ctx);
patterns.add<ConvOpLowering>(ctx);
patterns.add<ReluOpLowering>(ctx);
applyPartialConversion(module, target, std::move(patterns));
\end{lstlisting}

\subsection{패턴 6: 익명 namespace = Python \_private}
\begin{lstlisting}[language=C++]
namespace {
    struct ConvOpLowering : public OpConversionPattern<gawee::ConvOp> { ... };
} // namespace -- 이 파일에서만 접근 가능
\end{lstlisting}

\subsection{패턴 7: Lowering 패턴 뼈대}
\begin{lstlisting}[language=C++]
struct MyOpLowering : public OpConversionPattern<gawee::MyOp> {
    using OpConversionPattern::OpConversionPattern;
    LogicalResult
    matchAndRewrite(gawee::MyOp op, OpAdaptor adaptor,
                    ConversionPatternRewriter &rewriter) const override {
        Location loc = op.getLoc();
        Value input = adaptor.getInput();           // 1. 추출
        Value empty = tensor::EmptyOp::create(...); // 2. 생성
        rewriter.replaceOp(op, newResults);          // 3. 교체
        return success();
    }
};
\end{lstlisting}

\newpage

% ============================================================
\section{2c장. Modern C++ 코드 패턴}
% ============================================================

\subsection{RAII}
소멸자가 자원 해제. \texttt{with}문 필요 없음.
\begin{lstlisting}[language=C++]
{
    std::ifstream file("data.txt");  // 생성자에서 open
    // ...
}  // 소멸자에서 자동 close
\end{lstlisting}

\subsection{Rule of Zero}
멤버를 RAII 타입(\texttt{string}, \texttt{vector}, \texttt{unique\_ptr})으로만 쓰면 생성자/소멸자를 직접 안 써도 됨.

\subsection{Factory 함수}
\begin{lstlisting}[language=C++]
std::unique_ptr<Optimizer> createOptimizer(const Config &config) {
    if (config.type == "adam")
        return std::make_unique<Adam>(config.lr);
    return std::make_unique<SGD>(config.lr);
}
\end{lstlisting}

\subsection{enum class + switch}
\begin{lstlisting}[language=C++]
enum class Color { Red, Green, Blue };
switch (c) {
case Color::Red:   return "red";
case Color::Green: return "green";
case Color::Blue:  return "blue";
}
\end{lstlisting}

\subsection{const T\& 매개변수}
\begin{lstlisting}
sizeof(T) <= 16?   -> 값으로 (int, bool, float, pointer)
읽기만?            -> const T&
수정 필요?         -> T&
소유권 넘기기?     -> T&& 또는 unique_ptr<T>
\end{lstlisting}

\subsection{std::optional}
\begin{lstlisting}[language=C++]
std::optional<User> findUser(const std::string &name) {
    auto it = db.find(name);
    if (it != db.end()) return it->second;
    return std::nullopt;  // Python의 None
}
auto name = findName(id).value_or("unknown");
\end{lstlisting}

\subsection{Virtual Interface}
\begin{lstlisting}[language=C++]
class Serializer {
public:
    virtual ~Serializer() = default;       // 가상 소멸자 필수!
    virtual std::vector<uint8_t> serialize(const Data &d) = 0;
};
class JsonSerializer : public Serializer {
    std::vector<uint8_t> serialize(const Data &d) override { ... }
};
\end{lstlisting}

\subsection{CRTP}
\begin{lstlisting}[language=C++]
template <typename Derived>
class Printable {
public:
    void print() const {
        const auto &self = static_cast<const Derived &>(*this);
        std::cout << self.toString() << "\n";
    }
};
class Dog : public Printable<Dog> {
    std::string toString() const { return "Dog: woof"; }
};
\end{lstlisting}

\subsection{함정 모음}
\begin{longtable}{p{4cm}p{9cm}}
\toprule
함정 & 핵심 \\
\midrule
Dangling Reference & 로컬 변수의 참조/포인터를 반환하면 crash \\
Object Slicing & 자식을 부모 값으로 복사하면 자식 부분이 잘림 \\
초기화 순서 & 멤버는 선언 순서대로 초기화 (리스트 순서 아님) \\
Implicit Conversion & \texttt{explicit} 없으면 생성자가 몰래 타입 변환 \\
헤더/구현 분리 & 템플릿은 .h에, 나머지 구현은 .cpp에 \\
\bottomrule
\end{longtable}

\newpage

% ============================================================
\section{3장. C++ 디자인 패턴 --- MLIR에서 쓰이는 것만}
% ============================================================

\begin{longtable}{p{4cm}p{9cm}}
\toprule
패턴 & 핵심 \\
\midrule
Builder & \texttt{OpBuilder}가 상태(삽입 위치)를 관리하며 연산 생성 \\
Visitor & op\_type별 if-else 분기 \\
Registry & \texttt{patterns.add<T>(ctx)}로 등록, 프레임워크가 자동 매칭 \\
Template Method & 부모가 전체 구조, 자식은 \texttt{matchAndRewrite}만 구현 \\
\bottomrule
\end{longtable}

\newpage

% ============================================================
\section{4장. CRTP}
% ============================================================

\textbf{CRTP = 자기 자신의 타입을 부모에게 전달하는 패턴}

\begin{lstlisting}[language=C++]
template <typename Derived>
class Base {
  void execute() {
    static_cast<Derived*>(this)->runOnOperation(); // 컴파일 타임 결정
  }
};

struct MyPass : public PassWrapper<MyPass, OperationPass<ModuleOp>> {
  void runOnOperation() override { /* ... */ }
};
\end{lstlisting}

\textbf{읽기 규칙}: \texttt{struct A : public B<A, C>}
\begin{enumerate}[nosep]
\item A는 B를 상속
\item B는 A의 타입 정보를 컴파일 타임에 안다
\item C는 추가 설정
\end{enumerate}

\newpage

% ============================================================
\section{5장. MLIR 핵심 개념}
% ============================================================

\subsection{전체 구조}
\begin{lstlisting}
Gawee (고수준)  ->  Linalg (중간)  ->  LLVM (저수준)
gawee.conv         linalg.conv_2d     llvm.call
gawee.relu         linalg.generic     llvm.store
\end{lstlisting}

\subsection{Operation 구조}
\begin{lstlisting}
%result = gawee.conv %input, %weight, %bias {strides = [2, 2]}
          : tensor<...> -> tensor<...>

%result      = SSA Value (결과)
gawee.conv   = Op 이름
%input, ...  = Operands (런타임 값)
{strides}    = Attributes (컴파일 타임 상수)
tensor<...>  = Types
\end{lstlisting}

\subsection{Lowering 3요소}
\begin{enumerate}[nosep]
\item \textbf{ConversionTarget} --- 합법/불법 Op 정의
\item \textbf{ConversionPattern} --- matchAndRewrite 구현
\item \textbf{applyPartialConversion} --- 실행
\end{enumerate}

\subsection{matchAndRewrite 3단계}
\begin{enumerate}[nosep]
\item \texttt{adaptor}에서 입력값 추출
\item \texttt{rewriter}로 새 Op 생성
\item \texttt{rewriter.replaceOp}으로 교체
\end{enumerate}

\subsection{OpAdaptor}
\begin{itemize}[nosep]
\item \texttt{adaptor.getX()} --- 변환된 operand (변환 후)
\item \texttt{op.getXAttr()} --- 속성 (변환 불필요)
\end{itemize}

\subsection{Pass Pipeline}
\begin{lstlisting}
Gawee -> Linalg -> Bufferize -> SCF -> CF -> LLVM
고수준   선형대수   메모리화    루프  분기  기계어
\end{lstlisting}

\newpage

% ============================================================
\section{6장. TableGen}
% ============================================================

\subsection{전체 흐름}
\begin{lstlisting}
.td 파일 -> mlir-tblgen -> .inc 파일 -> .h/.cpp에서 #include
\end{lstlisting}

\subsection{Op 정의 예시}
\begin{lstlisting}
def Gawee_ConvOp : Gawee_Op<"conv", []> {
  let arguments = (ins
    AnyTensor:$input,               // Operand
    AnyTensor:$weight,              // Operand
    DenseI64ArrayAttr:$strides,     // Attribute
    DenseI64ArrayAttr:$padding      // Attribute
  );
  let results = (outs AnyTensor:$output);
}
\end{lstlisting}

\subsection{자동 생성되는 것}
C++ 클래스, 접근자 메서드, 파서, 프린터, Verifier

\subsection{\$name $\to$ C++ 접근자}
\begin{itemize}[nosep]
\item \texttt{\$input} $\to$ \texttt{getInput()}
\item \texttt{\$strides} $\to$ \texttt{getStrides()} / \texttt{getStridesAttr()}
\end{itemize}

\newpage

% ============================================================
\section{7장. CMake 빌드 시스템}
% ============================================================

\subsection{CMakeLists.txt 구조}
\begin{lstlisting}[language=cmake]
cmake_minimum_required(VERSION 3.20)
project(GaweeMLIR LANGUAGES C CXX)
set(CMAKE_CXX_STANDARD 17)

# RTTI 설정
if(NOT LLVM_ENABLE_RTTI)
  add_compile_options(-fno-rtti)
endif()

# MLIR 찾기
find_package(MLIR REQUIRED CONFIG)
include(TableGen)
include(AddLLVM)
include(AddMLIR)

# 라이브러리
add_library(GaweeDialect lib/Gawee/GaweeDialect.cpp)
target_link_libraries(GaweeDialect PUBLIC MLIRIR MLIRSupport)

# 실행파일
add_executable(gawee-opt tools/gawee-opt.cpp)
target_link_libraries(gawee-opt PRIVATE GaweeDialect GaweeConversion)
\end{lstlisting}

\subsection{빈출 에러}
\begin{longtable}{p{6cm}p{7cm}}
\toprule
에러 & 해결 \\
\midrule
\texttt{typeinfo for mlir::Pass} & RTTI 불일치 $\to$ \texttt{-fno-rtti} \\
\texttt{undefined reference} & 라이브러리 미링크 \\
\texttt{file not found} & include 경로 누락 \\
\texttt{package not found} & \texttt{MLIR\_DIR} 미설정 \\
\bottomrule
\end{longtable}

\newpage

% ============================================================
\section{Mini KV Store 설계 명세서 (SPEC)}
% ============================================================

\subsection{커맨드}
\begin{longtable}{p{4cm}p{9cm}}
\toprule
커맨드 & 설명 \\
\midrule
\texttt{set <key> <value>} & 키-값 쌍 저장 \\
\texttt{get <key>} & 값 조회 \\
\texttt{delete <key>} & 삭제 \\
\texttt{list} & 모든 키-값 알파벳 순 출력 \\
\texttt{count} & 저장된 키 개수 \\
\texttt{export <format>} & json 또는 csv 출력 \\
\bottomrule
\end{longtable}

\subsection{클래스 설계}
\begin{itemize}[nosep]
\item \textbf{KVStore} --- 데이터 저장 + 파일 I/O (\texttt{map}, \texttt{optional}, RAII)
\item \textbf{Command} --- 추상 인터페이스 + 구체 커맨드 (Virtual Interface, Factory)
\item \textbf{Exporter} --- 출력 포맷 다형성 (Strategy Pattern)
\end{itemize}

\newpage

% ============================================================
\section{진행 현황 (PROGRESS)}
% ============================================================

\begin{longtable}{p{1cm}p{5cm}p{6cm}}
\toprule
단계 & 만드는 것 & 핵심 개념 \\
\midrule
1 & 디렉토리 + 빈 파일 & 프로젝트 구조 관례 \\
2 & CMakeLists.txt + 첫 빌드 & CMake 기본 명령어 \\
3 & KVStore 헤더 선언 & 클래스 선언, const \&, explicit \\
4 & KVStore 구현 (메모리) & std::map, std::optional \\
5 & save() & std::ofstream, RAII \\
6 & load() & std::ifstream, 문자열 파싱 \\
7 & Command 인터페이스 선언 & 순수 가상 함수, 가상 소멸자 \\
8 & Set/Get 커맨드 구현 & override, 가상 함수 호출 \\
9 & 나머지 커맨드 & 패턴 반복 연습 \\
10 & parseCommand + main & Factory, unique\_ptr, argc/argv \\
11 & Exporter + CSV & ostream \&, Strategy \\
12 & JSON + 팩토리 & Factory 두 번째 연습 \\
13 & 에러 처리 + test.sh & 방어적 프로그래밍 \\
\bottomrule
\end{longtable}

전체 약 6.5시간 (13단계 $\times$ 30분)

\end{document}
